\chapter{How to Read the Line Atlas}
\label{chap:foundations}

Spectroscopy is a translation problem.  The detector records counts as a
function of coordinate; the astronomer wants temperature, density, chemical
composition, radiation field, velocity, geometry, and evolutionary history.
A spectral transition connects those descriptions, but never by itself.  Its
meaning depends on how the relevant state was populated, how photons propagated,
which material shared the resolution element, and how the instrument transformed
the incoming radiation.  This opening chapter establishes the conventions used
throughout the atlas and a disciplined way to move from line identification to
physical inference.

\begin{conceptbox}[The central rule]
A line is evidence for a transition of a species.  Calling it a tracer of a
temperature, phase, or process is a model-dependent inference.  Preserve the
distinction between what the spectrum measures and what a physical model adds.
\end{conceptbox}

\section{The information carried by a line}

An isolated feature may provide several nearly independent observables:
\begin{enumerate}[label=\arabic*.]
  \item its centroid identifies a candidate transition and line-of-sight velocity;
  \item its integrated flux or equivalent width measures strength relative to a
        calibrated continuum;
  \item its width and profile encode thermal motion, turbulence, rotation,
        outflow, opacity, unresolved components, and the instrumental response;
  \item a resolved doublet or multiplet tests saturation, density, optical depth,
        and line identification;
  \item ratios to other transitions constrain excitation, ionization, attenuation,
        abundance, and the shape of the radiation field.
\end{enumerate}
These observables do not all imply the same spatial weighting.  Emissivity often
scales as a product of two densities, while absorption optical depth is linear in
the population of the lower state.  Bright emission therefore favors dense
structures; absorption can detect diffuse material along a long path.

For a transition between levels $u$ and $l$,
\begin{equation}
  h\nu_0 = E_u-E_l, \qquad \lambda_0 = \frac{c}{\nu_0}.
  \label{eq:foundation-transition}
\end{equation}
The rest coordinate is a laboratory property, whereas the observed coordinate is
altered by cosmological redshift and peculiar velocity.  To first order in
$v/c$, a small velocity displacement from a systemic wavelength obeys
\begin{equation}
  \Delta v \simeq c\,\frac{\lambda-\lambda_{\rm sys}}
                              {\lambda_{\rm sys}}.
\end{equation}
For high precision or large velocities one must state and use a relativistic
Doppler convention rather than mixing radio, optical, and relativistic velocity
definitions.

\begin{figure}[tbp]
\centering
\begin{tikzpicture}[x=1.2cm,y=1.05cm,>=Latex]
  \foreach \y/\lab in {0/$E_1$,2/$E_2$,4/$E_3$}{
    \draw[AtlasNavy,line width=1pt] (0,\y)--(4.2,\y) node[right]{\lab};
  }
  \draw[->,AtlasCyan,line width=1.4pt] (0.9,0.08)--(0.9,1.92)
    node[midway,left,align=right,font=\small]{absorption\\$B_{12}J_\nu$};
  \draw[->,AtlasCoral,line width=1.4pt] (2.25,1.92)--(2.25,0.08)
    node[midway,left,align=right,font=\small]{spontaneous\\$A_{21}$};
  \draw[->,AtlasAmber,line width=1.4pt] (3.6,1.92)--(3.6,0.08)
    node[midway,right,align=left,font=\small]{stimulated\\$B_{21}J_\nu$};
  \draw[->,AtlasCyan,dashed,line width=1.2pt] (0.9,2.08)--(0.9,3.92)
    node[midway,left,align=right,font=\small]{collisional\\excitation $C_{23}$};
  \node[anchor=west,text width=4.8cm,align=left] at (5.0,3.4) {The same observed transition can be populated by radiative or collisional routes.};
  \node[anchor=west,text width=4.8cm,align=left] at (5.0,1.2) {Statistical equilibrium balances every route into and out of each level; a line is not assigned a unique formation condition by wavelength alone.};
\end{tikzpicture}
\caption{The principal bound--bound population and depopulation routes. The Einstein coefficients describe radiative transitions; collisional rates depend on collider density and the velocity-averaged cross section.}
\label{fig:line-processes}
\end{figure}


\section{Transition families}

The atlas groups lines by mechanism as well as wavelength.  The mechanism
determines the relevant selection rules, characteristic lifetime, and sensitivity
to density or radiation.

\begin{table}[tbp]
\centering
\caption{A physical taxonomy of common astronomical transitions.  The final
column gives a tendency, not an exclusive rule.}
\label{tab:transition-families}
\small
\begin{tabularx}{\textwidth}{@{}lXX X@{}}
\toprule
Family & Physical origin & Typical examples & Often constrains \\
\midrule
Permitted resonance & Electric-dipole transition from a populated lower level &
Ly$\alpha$, C\,IV, Mg\,II & column density, covering fraction, kinematics \\
Recombination & Cascades following radiative recombination &
H$\alpha$, H$\beta$, He\,II, radio recombination lines & ionizing-photon rate,
attenuation, emission measure \\
Forbidden or semi-forbidden & Long-lived state populated collisionally &
[O\,III], [S\,II], C\,III] & temperature, density, abundance, excitation source \\
Fine structure & Splitting within a ground electronic term &
[C\,II] $158\,\mu$m, [O\,I] $63\,\mu$m & gas cooling, collider density, PDR or
ionized-gas conditions \\
Rotational or inversion & Molecular rotation or tunnelling between configurations &
CO, HCN, NH$_3$ & molecular excitation, mass proxies, temperature, dense gas \\
Hyperfine & Coupling of nuclear and electronic angular momentum &
H\,I 21 cm, OH 18 cm & neutral gas, spin temperature, magnetic field \\
Inner-shell or nuclear & Atomic K-shell rearrangement or nuclear decay &
Fe K$\alpha$, 511 keV, $^{26}$Al & irradiation, compact objects, nucleosynthesis \\
\bottomrule
\end{tabularx}
\end{table}

The adjective ``forbidden'' means forbidden by the leading electric-dipole
selection rules, not impossible.  In a sufficiently low-density plasma the upper
state survives long enough to radiate.  At high collider density, de-excitation
may occur collisionally before a photon escapes; this is the origin of critical
density diagnostics.  Likewise, a molecular line with a large dipole moment is
not automatically a unique dense-gas tracer: radiative trapping, subthermal
excitation, abundance, infrared pumping, and beam dilution all matter.

\section{Reading species and transition labels}

Astronomical ion notation counts the neutral atom as stage I, so O\,VI is
O$^{5+}$ rather than O$^{6+}$.  Square brackets denote a forbidden transition,
as in [O\,III] $\lambda5007$; a closing bracket alone conventionally marks a
semi-forbidden transition, as in C\,III].  A wavelength following $\lambda$ is
usually in Angstrom unless another unit is shown.  Molecular notation instead
specifies rotational, vibrational, inversion, or hyperfine quantum numbers, for
example CO($J=1\!\rightarrow\!0$) or H$_2$ 1--0 S(1).

The ionization threshold needed to create an ion is not the energy of the listed
line.  The former describes removal of an electron; the latter is the separation
of two bound levels.  Confusing them can lead to incorrect claims about the
hardness of the radiation required to produce a feature.  Ion fractions also
depend on recombination and charge exchange, not only threshold energy
\citep{OsterbrockFerland2006,Draine2011}.

\section{Coordinates, units, and conventions}

\subsection{Vacuum and standard air wavelengths}

This edition stores atomic wavelengths in vacuum unless an entry explicitly says
otherwise.  Familiar ground-based optical names often use standard air values:
H$\alpha$ is commonly called $6562.8$~\AA\ in air but is approximately
$6564.6$~\AA\ in vacuum.  The difference is scientifically significant for
precision velocities.  Conversion depends on a refractive-index prescription and
its validity range, so an analysis must record both the input convention and the
formula used.  NIST distinguishes observed and Ritz wavelengths and documents
the medium used by its tabulations \citep{nist-asd}.

\subsection{Use the native coordinate}

At centimetre and millimetre wavelengths, observers tune to a rest frequency and
should retain it rather than relying on a rounded wavelength conversion.  X-ray
features are often reported in energy, while optical and ultraviolet features are
normally reported in wavelength.  The equivalent representations obey
\begin{equation}
  E = h\nu = \frac{hc}{\lambda},
  \qquad
  \left|\frac{\Delta E}{E}\right|
  =\left|\frac{\Delta\nu}{\nu}\right|
  =\left|\frac{\Delta\lambda}{\lambda}\right|
\end{equation}
for small intervals.  Significant digits should follow the authoritative
laboratory uncertainty and component definition, not the numerical precision of
a unit conversion.

\subsection{Flux-density transformations}

Flux density per unit frequency and per unit wavelength are related by a Jacobian,
\begin{equation}
 f_\lambda = f_\nu\left|\frac{\mathrm d\nu}{\mathrm d\lambda}\right|
 =f_\nu\frac{c}{\lambda^2},
 \qquad \lambda f_\lambda=\nu f_\nu.
\end{equation}
A spectrum that is flat in $f_\nu$ is therefore not flat in $f_\lambda$.
Integrated line flux remains an energy flux after integrating over a consistently
defined coordinate, but equivalent widths and continuum slopes must retain their
coordinate convention \citep{RybickiLightman1979}.

\section{From identification to inference}

A credible line identification is a network of consistency checks rather than a
nearest-wavelength match.  The following sequence is used throughout this book.

\begin{enumerate}[label=\textbf{Step \arabic*:},leftmargin=*]
  \item \textbf{Calibrate the coordinate.} Establish the wavelength solution,
  barycentric or local-standard-of-rest correction, air/vacuum convention, and
  uncertainty.
  \item \textbf{Generate candidate transitions.} Use the allowed redshift range,
  spectral coverage, expected source class, and a current laboratory database.
  \item \textbf{Demand companions.} Search for doublet partners, other members of
  the same series, adjacent ionization stages, and physically expected continuum
  features.
  \item \textbf{Test the profile.} Compare component velocities and widths after
  applying the line-spread function; inspect blends, sky residuals, detector
  defects, and telluric absorption.
  \item \textbf{Separate measurement from model.} Report flux, equivalent width,
  centroid, width, covariance, and limits before translating them into abundance,
  density, temperature, or mass.
  \item \textbf{Stress-test alternatives.} Vary continuum placement, component
  count, attenuation law, atomic data, ionization mechanism, geometry, and priors.
\end{enumerate}

\begin{cautionbox}[A non-detection is not zero]
A non-detection constrains a line only through the noise model, assumed profile,
searched velocity interval, and selection procedure.  Report an upper limit with
its confidence convention and assumed width; preserve the likelihood or sensitivity
curve when it will enter population inference.
\end{cautionbox}

\section{How each atlas entry is constructed}

Each generated entry records an operational rest coordinate, physical mechanism,
common environments, useful tracers, an indicative temperature regime, density
guidance, diagnostic companions, observing cautions, and source keys.  The source
hierarchy is deliberate: critically evaluated atomic or molecular databases supply
rest data; primary papers and reviews supply interpretation; observatory documents
supply current instrumental limits.  The prose is a guide to discovery, not a
replacement for a live database query or a forward model.

The phrase ``temperature regime'' is especially easy to overread.  A quoted
collisional-ionization-equilibrium peak is an orientation point.  Photoionization,
non-equilibrium cooling, shocks, mixing, charge exchange, optical depth, and an
unresolved multiphase sightline can produce the same ion under different kinetic
conditions.  Similarly, a critical density depends on the collider and temperature;
an effective excitation density also depends on opacity and column.

\section{A compact reporting standard}

For a result to be reproducible, a line measurement should state at least:
\begin{itemize}
  \item transition and component, adopted rest coordinate, source database and
        version, and air/vacuum convention;
  \item reference frame, redshift or velocity convention, and systemic-redshift
        uncertainty;
  \item instrument, configuration, resolving power, sampling, aperture, and
        line-spread-function treatment;
  \item continuum model, integration or fitting window, blend model, and
        attenuation correction;
  \item flux, equivalent width, centroid, width, uncertainties and covariance, or
        the complete upper-limit prescription;
  \item physical model, radiation field, abundance pattern, atomic or molecular
        data, priors, and sensitivity tests used for interpretation.
\end{itemize}

This standard may look longer than a table cell, but most fields belong in a
machine-readable companion product.  The purpose of the atlas is to make those
dependencies visible early, before a memorable line name hardens into an
overconfident physical claim.
